\documentclass[11pt]{article}
\usepackage{amsmath,amssymb,amsthm,hyperref}
\newtheorem{proposition}{Proposition}
\title{Arithmetic Storage and the Cost of a Flat Decoder}
\author{}
\date{November 23, 2026}
\begin{document}
\maketitle

\paragraph{Scope.}
The statements below audit an arithmetic construction, not the requested
upper bound. In particular, a fixed decoder is an additional restriction.
The manuscript allows all head parameters to depend on the target.

\begin{proposition}[One positive affine integer stores every table]
Let $n\ge1$, $N=2^n$, and
$t(x)=\sum_{i=0}^{n-1}2^i x_i$ on $\{0,1\}^n$.
There is a positive integer $M$, depending only on $n$, such that every
subset $S$ of the cube has an integer $T$ with $1\le T\le M$ satisfying
\[
 x\in S\quad\Longleftrightarrow\quad\gcd(T+t(x),M)>1.
\]
The affine form $T+t(x)$ is positive and all its variable coefficients
are strictly positive.
\end{proposition}

\begin{proof}
Choose distinct primes $p_v>2N$ for the $N$ vertices $v$ and put
$M=\prod_v p_v$. By the Chinese remainder theorem, choose the unique
representative $T\in\{1,\ldots,M\}$ satisfying
\[
 T\equiv
 \begin{cases}
 -t(v)&v\in S,\\
 N-t(v)&v\notin S
 \end{cases}
 \pmod {p_v}.
\]
If $v\in S$, then $T+t(x)\equiv t(x)-t(v)\pmod {p_v}$.
The difference lies between $-(N-1)$ and $N-1$, so divisibility occurs
exactly when $x=v$. If $v\notin S$, the corresponding residue is
$N+t(x)-t(v)$, an integer between $1$ and $2N-1$, and divisibility
never occurs. A common prime divisor with $M$ therefore exists exactly
at the vertices of $S$. Positivity of $T$ proves the remaining claims.
\end{proof}

The greatest-common-divisor test in this proposition is a decoding
operation. It is not an attention head or a linear readout. The next
statement quantifies one attempted replacement for that operation.

\begin{proposition}[Affine dependence on one stored real parameter]
Let $X$ have $N$ elements. Fix a family
\[
 F_T(x)=c+\sum_{h=1}^H
 \frac{a_{h,x}T+b_{h,x}}{u_{h,x}T+v_{h,x}},
\]
where all displayed coefficients are independent of $T$.
Suppose there is a set $E\subset\mathbb R$ such that every denominator
is positive for every $(T,x)\in E\times X$, all scores are nonzero
there, and the vectors $\bigl(\operatorname{sgn}F_T(x)\bigr)_{x\in X}$,
as $T$ ranges over $E$, include all $2^N$ sign vectors. Then
\[
 NH\ge 2^N-1.
\]
If the common bias is affine in $T$, the same proof gives
$N(H+1)\ge2^N-1$.
\end{proposition}

\begin{proof}
Choose one parameter for each distinct sign vector and sort these
$2^N$ real numbers. They are distinct because a fixed parameter has
only one sign vector. For each $x$, clear the $H$ denominators to obtain
a polynomial $P_x(T)$ of degree at most $H$. It is not identically
zero, since every selected score is nonzero. At each selected
parameter its sign equals the sign of $F_T(x)$.

Every consecutive pair in the sorted list has different sign vectors.
For at least one $x$, the corresponding values of $P_x$ have opposite
signs. The intermediate value theorem gives a root in that open
parameter interval. The intervals are disjoint. Charging each
interval to one such polynomial gives at least $2^N-1$ distinct
polynomial roots counted with their polynomial labels. On the other
hand, the $N$ nonzero polynomials have at most $NH$ real roots in
total. This proves the inequality. An affine bias increases the
degree bound by at most one.
\end{proof}

The coefficients in the second proposition need not be affine in
$x$, so its obstruction also applies to that more generous decoder.
It excludes a fixed small affine-fraction device that reads the
Chinese-remainder code merely by inserting $T$ affinely into its
coefficients. It does not exclude target-dependent denominators with
arbitrary dependence on the target, as required in the problem.

\paragraph{Differential encodings.}
Sawczuk and Gnang's \emph{Symbolic Listings as Computation},
\href{https://arxiv.org/pdf/2402.15885}{arXiv:2402.15885},
Definition 2.1 and Section 2, store a Boolean table in polynomial
coefficients and retrieve entries with input-selected differential
operators. Gnang and Xu's
\emph{On Partial Differential Encodings of Boolean Functions},
\href{https://arxiv.org/pdf/2008.06801}{arXiv:2008.06801},
Proposition 1 and Definition 2, similarly distinguish a polynomial
encoding from coefficient retrieval. The complex factorization in
Proposition 1 has up to $2^n-1$ affine factors. Neither the
differentiation operation nor the complex factors are admissible
affine fractions in the manuscript. A small Chow rank therefore
does not itself supply a head-count bound.

\paragraph{Status of the mechanism.}
Arithmetic storage is explicit, but no decoder with the required
head count has been constructed. Replacing the decoder by an
unproved universal sign-realization lemma would restate the target
gap. This route remains blocked until an admissible decoding
construction supplies the missing quantitative bound.

\end{document}
